% Generated draft: requires mathematical and source review.
\documentclass[11pt]{article}
\usepackage[margin=25mm]{geometry}
\usepackage{fontspec}
\setmainfont{Latin Modern Roman}
\usepackage{amsmath,amssymb,mathtools,unicode-math}
\setmathfont{Latin Modern Math}
\usepackage{enumitem,needspace,xurl,hyperref}
\setlength{\parindent}{0pt}
\setlength{\parskip}{5pt}
\setlength{\emergencystretch}{3em}
\newcommand{\N}{\mathbb{N}}
\newcommand{\Z}{\mathbb{Z}}
\newcommand{\Q}{\mathbb{Q}}
\newcommand{\R}{\mathbb{R}}
\newcommand{\C}{\mathbb{C}}
\newcommand{\F}{\mathbb{F}}
\newcommand{\A}{\mathbb{A}}
\newcommand{\PP}{\mathbb{P}}
\newcommand{\E}{\mathbb{E}}
\newcommand{\eps}{\varepsilon}
\newcommand{\dd}{\,\mathrm{d}}
\begin{document}
\section*{UnsolvedMath ID: 2301038}
\textbf{Source problem code:} catalogue entry
\subsection{Definitions and mathematical statement}
Let $\mathbb D=\{z\in\C:|z|<1\}$ and let $T(r,f)$ denote the Nevanlinna characteristic of a meromorphic function in $\mathbb D$. For a nonconstant such function, set
\[
 \alpha(f)=\limsup_{r\uparrow1}\frac{T(r,f)}{-\log(1-r)}<\infty.
\]
For a meromorphic expression $\Psi$, the required conclusion is
\[
 \forall a\in\C\setminus\{0\},\qquad
 \#\{z\in\mathbb D:\Psi(z)=a\}=\infty.
\]
Thus zero is the only permitted exceptional finite value, and infinitely many distinct preimages are required.

\textbf{First question.} With the usual differential-monomial convention $k\geq0$ and nonnegative integral exponents $m_0,\ldots,m_k$, put
\[
 \Psi=f^{m_0}(f')^{m_1}\cdots(f^{(k)})^{m_k}.
\]
The source reports the conclusion above when $f$ is meromorphic, $m_0\geq3$, and $\alpha(f)>2/(m_0-2)$; for analytic $f$, it reports the improved hypothesis $m_0\geq2$ and $\alpha(f)>2/(m_0-1)$. For which smaller values of $\alpha(f)$ does the same conclusion hold, retaining the corresponding function class and exponent assumptions? In particular, the respective remaining ranges include the threshold endpoints and values below them; no optimal threshold is supplied.

\textbf{Second question.} Let $f$ additionally have only finitely many zeros and poles in $\mathbb D$, let $l\geq1$ be an integer, and put
\[
 \Psi=\sum_{\nu=0}^{l}a_\nu f^{(\nu)},\qquad f^{(0)}=f,
 \qquad T(r,a_\nu)=O(T(r,f))\quad(r\uparrow1)
\]
for every coefficient $a_\nu$. The source calls these coefficients functions in $\mathbb D$; their meromorphic character is the natural convention needed to interpret $T$. It reports that a nonconstant $\Psi$ has the required infinite-preimage property if
\[
 \alpha(f)>\tfrac12l(l+1)+1.
\]
For which smaller $\alpha(f)$ does the conclusion remain valid under the stated hypotheses?

The large-$O$ coefficient condition is transcribed exactly. Arbitrary coefficients of size comparable to $f$ can permit cancellation, so the recorded theorem needs source checking: for example, the allowed shape $a_0=g/f$, $a_\nu=0$ for $\nu\geq1$ gives $\Psi=g$ whenever its characteristic condition holds. The record gives no restriction ruling out this dependence. A small-$o$ condition, constant coefficients, or other added restrictions are not substituted.

Both questions are attributed to L. R. Sons. The supplied summary says they were open in Hayman's 2018 edition and no recent resolution was located. The quoted sufficient conditions and status are source reports, not independently verified theorems.
\subsection{Short English statement}
Determine how far the given logarithmic-growth thresholds can be lowered while differential monomials, and separately linear differential expressions, still take every nonzero finite value infinitely often in the unit disc. The second part's large-$O$ coefficient condition needs checking because it permits cancellation.
\subsection{Sources}
\begin{itemize}
\item Local UnsolvedMath registry, permanent ID 2301038, source code catalogue entry, ``Research Problems in Function Theory --- Problem 1.38''; attribution: L. R. Sons. Supplied status open, classification OPEN-TRIAGE, updated 2026-08-06.
\item Hayman, \emph{Research Problems in Function Theory} (2018), Problem 1.38, \url{https://arxiv.org/abs/1809.07200}. The record reports source-TeX extraction and access on 2026-07-29; the 2018 update was retained for status triage. Current status and rights require review; the record describes publicly available arXiv source TeX.
\item The sufficient bounds and no-recent-resolution-found assessment are supplied in this record without separate papers or theorem references.
\end{itemize}
\end{document}
